% system_overview.tex -- Methodology: system and platform (draft 4 Oct 2026).
% Restructured 9 Oct (structure audit): the platform paragraphs and fig:geometry moved here from
% hardware.tex, so the geometry is defined before the model; the duplicate reference and takeoff
% paragraphs cut (the interpolation sentence moved to the planner's output).

\subsection{System and platform}
\label{sec:overview}

\subsubsection{Nodes and safety}

All control runs centrally, on one ground computer under ROS~2 (Robot Operating System~2). Motion capture (OptiTrack) is the
only sensing. It gives the pose of each UAV and of the payload, and nothing on the payload is
instrumented. Each UAV's flight controller runs Betaflight in rate mode. It receives a collective
throttle and three body-rate stick commands over its own ExpressLRS (ELRS) radio link, and closes only its
inner rate loop. Figure~\ref{fig:nodes} shows the nodes and the information they exchange.

\begin{figure}[h]
\centering
\begin{tikzpicture}[>=Stealth,
    blk/.style={draw, thick, rounded corners=2pt, align=center, minimum height=0.95cm, font=\small},
    lab/.style={font=\scriptsize, align=center}]
  \node[blk, minimum width=2.4cm] (mocap) at (0,0) {Motion capture\\ \scriptsize poses};
  \node[blk, minimum width=2.7cm] (plan) at (4.9,0) {Load--cable planner\\ \scriptsize 10\,Hz};
  \node[blk, minimum width=2.4cm] (fm) at (9.5,0) {Fleet manager\\ \scriptsize phases, safety};
  \node[blk, minimum width=1.9cm] (op) at (13.4,0) {Operator};
  \node[blk, minimum width=2.7cm] (trk) at (4.9,-2.2) {Tracker $i$\\ \scriptsize 50\,Hz, one per UAV};
  \node[blk, minimum width=2.7cm] (uav) at (4.9,-4.4) {UAV $i$\\ \scriptsize ELRS link, Betaflight rate loop};
  \draw[->, thick] (mocap) -- node[lab, above] {payload\\ pose} (plan);
  \draw[->, thick] (mocap.south) |- node[lab, pos=0.75, above] {UAV pose} (trk.west);
  \draw[->, thick] (plan) -- node[lab, right] {horizon for UAV $i$} (trk);
  \draw[->, thick] (trk) -- node[lab, right] {throttle, rate sticks, magnet} (uav);
  \draw[->, thick] (fm) -- node[lab, above] {phase} (plan);
  \draw[->, thick] (op) -- node[lab, above] {commands} (fm);
  \draw[->, thick] ([xshift=-0.35cm]fm.south) |- node[lab, pos=0.75, above] {arm, disarm} ([yshift=0.18cm]trk.east);
  \draw[->, thick, dashed] ([yshift=-0.18cm]trk.east) -| node[lab, pos=0.25, below] {fault} ([xshift=0.35cm]fm.south);
\end{tikzpicture}
\caption[Nodes and links]{Nodes and links. Solid: commands, references and poses; dashed: a tracker
reporting a fault, which disarms the whole fleet. The operator's commands are ARM, TAKEOFF, LAND,
DETACH and ATTACH.}
\label{fig:nodes}
\end{figure}

The system has three layers.
\begin{description}
    \item[Planner.] The centralised load--cable MPC of Section~\ref{sec:controller}. At 10\,Hz it
    plans the motion of the payload and of every rod, and publishes to each UAV its reference over
    the horizon: position, velocity, acceleration, the planned cable acceleration and the throttle
    feedforward. A second reference generator, the dissipative network of
    Section~\ref{sec:dissipative}, can replace it in flight; both publish the same reference format,
    so the trackers do not know which one is active.
    \item[Trackers.] One tracking MPC per UAV at 50\,Hz. Each follows its UAV's part of the plan and
    turns it into the radio commands.
    \item[Fleet manager.] It sequences the phases of a flight: ARM, TAKEOFF (the climb, the rod sweep
    to the working elevation and the lift), the trajectory, LAND, and the reconfiguration events
    DETACH and ATTACH (Section~\ref{sec:reconfig}).
\end{description}
\paragraph{Safety.} Every tracker checks its own UAV against a flight envelope: tilt, payload tilt
(both 60$^\circ$) and speed (3\,m/s). It also checks that its motion-capture pose is fresh (0.25\,s)
and that the planner's reference is fresh (1\,s). Any violation disarms the whole fleet, not the one
UAV: a single UAV losing thrust on a shared rigid payload drags the others with it. There is no
geofence; the operator holds a kill switch. On the rig, each UAV must also confirm that its flight
controller is armed before the fleet takes off. The planner stops iterating at a 60\,ms budget,
and a plan that has not converged by then is replaced by the previous one, shifted by a tick.
A UAV that disarms between ARM and TAKEOFF grounds the whole fleet and TAKEOFF is refused,
since three UAVs lifting a ring sized for four would drag it.

\subsubsection{Rig hardware}
\label{sec:hardware}

\paragraph{UAVs.} Four UAVs run Betaflight in rate mode on six-cell batteries. Each has its own
ELRS receiver and a ground transmitter on a USB serial port, so every UAV is commanded on its
own link at 100\,Hz. The throttle, three rate sticks, the arm switch and an auxiliary channel are
sent; the auxiliary channel switches the UAV's magnet and is held on from start-up, so a quiet link
cannot drop the payload. The flight controller's telemetry (battery voltage, flight mode, link
quality) comes back on the same link. As flown, each UAV weighs 0.55\,kg: an airframe and battery of
0.475\,kg, plus its rod and magnet.

\paragraph{Rods and payload.} Each UAV carries a rigid rod of 0.55\,m and 0.075\,kg with a ball joint
at both ends. The upper joint sits 4\,cm below the UAV's centre. The payload is a 0.86\,kg ring of
0.51\,m diameter and 40\,mm thickness, with twelve magnet plates every 30$^\circ$ on its top face at a
radius of 0.225\,m.
Each rod ends in an electromagnet that holds on a plate with 10--12\,N and releases under a jerk.
A layout is named by its plate numbers, anticlockwise from plate 0 on the ring's $+x$ axis
(Figure~\ref{fig:geometry}).

% fig_geometry.tex -- M3: payload and rod geometry (TikZ). Dimensions from the rig (30 Sep 2026):
% plates at r 0.225 m on a 500 mm ring, 12 plates every 30 deg anticlockwise from +x, rod 0.55 m,
% pivot 4 cm below the UAV centre, working elevation 45 deg. Drawn to scale (1 m = 9 cm).
\begin{figure}[h]
\centering
\begin{tikzpicture}[x=9cm, y=9cm, >=Stealth, font=\small]
  % ---------------- (a) top view ----------------
  \begin{scope}
    \fill[gray!12] (0,0) circle (0.25);
    \fill[white] (0,0) circle (0.19);
    \draw[gray!60, thin] (0,0) circle (0.25) (0,0) circle (0.19);
    % the even triangle left after detaching plate 3
    \draw[densely dashed, gray!70] (30:0.225) -- (150:0.225) -- (270:0.225) -- cycle;
    % plates
    \foreach \k in {0,...,11} {
      \pgfmathsetmacro{\a}{30*\k}
      \node[draw, fill=white, minimum size=0.11cm, inner sep=0pt, rotate=\a] at (\a:0.225) {};
      \node[font=\scriptsize, text=gray!70!black] at (\a:0.29) {\k};
    }
    % the four carrying plates
    \foreach \k in {1,3,5,9} { \fill[black] (30*\k:0.225) circle (0.012); }
    \draw[thick, black] (90:0.225) circle (0.022);
    % centre, axes, rho
    \fill (0,0) circle (0.008);
    \draw[->, thin, gray] (0,0) -- (0.09,0) node[below, font=\scriptsize, text=black] {$x_L$};
    \draw[->, thin, gray] (0,0) -- (0,0.09) node[left, font=\scriptsize, text=black] {$y_L$};
    \draw[->, thick] (0,0) -- (30:0.213) node[pos=0.62, above left, inner sep=1pt] {$\boldsymbol{\rho}_i$};
    \node[font=\scriptsize, align=left, anchor=north] at (0,-0.33) {$\bullet$ carrying plates 1, 3, 5, 9\quad $\bigcirc$ plate 3, detached};
    \node[anchor=north west] at (-0.36,0.36) {(a)};
  \end{scope}
  % ---------------- (b) side view of one rod (A plate, P pivot, C UAV centre) ----------------
  \begin{scope}[shift={(0.62,-0.30)}]
    \fill[gray!12] (-0.06,0) rectangle (0.0,0.03);
    \draw[gray!60, thin] (-0.06,0) rectangle (0.0,0.03);
    \fill[black] (-0.012,0.03) rectangle (0.012,0.038);
    \coordinate (A) at (0,0.034);
    \coordinate (P) at (0.3889,0.4229);
    \coordinate (C) at (0.3889,0.4629);
    \draw[very thick] (A) -- (P);
    \fill (A) circle (0.007) (P) circle (0.007);
    \draw[thick] (0.2989,0.4629) -- (0.4789,0.4629);
    \draw[gray] (0.2989,0.4689) ellipse (0.035 and 0.006) (0.4789,0.4689) ellipse (0.035 and 0.006);
    \fill (C) circle (0.006);
    \draw[->, thick] (A) -- (0.1131,0.1471) node[pos=0.55, above left, inner sep=1pt] {$\mathbf{q}_i$};
    \draw[gray] (0.16,0.034) arc (0:45:0.16);
    \draw[gray, thin] (A) -- (0.22,0.034);
    \node[font=\scriptsize] at (0.185,0.1105) {45$^\circ$};
    \node[font=\scriptsize, anchor=north west] at (0.2045,0.2185) {$l_i = 0.55$\,m};
    \draw[gray, thin, densely dotted] (P) -- (0.5689,0.4229) (C) -- (0.5689,0.4629);
    \draw[<->, thin] (0.5689,0.4229) -- (0.5689,0.4629) node[midway, right, font=\scriptsize] {$\boldsymbol{\delta}$: 4\,cm};
    \node[font=\scriptsize, anchor=north] at (-0.03,-0.005) {ring};
    \node[font=\scriptsize, anchor=south] at (0.3889,0.49) {UAV $i$};
    \node[anchor=north west] at (-0.08,0.66) {(b)};
  \end{scope}
\end{tikzpicture}
\caption[Payload and rod geometry]{Payload and rod geometry, to scale. (a) Top view of the 0.86\,kg ring
with its twelve magnet plates at a radius of 0.225\,m, numbered anticlockwise from the $x_L$ axis. The
four-UAV layout uses plates 1, 3, 5 and 9; detaching the UAV on plate 3 leaves the even triangle
(dashed), under which the ring hangs level. (b) Side view of one rod at the working elevation: the rod
direction $\mathbf{q}_i$ points from the plate to the UAV, and the rod pivots $\boldsymbol{\delta}$ = 4\,cm
below the UAV's centre.}
\label{fig:geometry}
\end{figure}


Figure~\ref{fig:rig_photo} shows the rig.

\pendingfig{fig:rig_photo}{The rig: four UAVs on rods, the ring with its twelve magnet plates, the motion-capture markers and the newcomer's magnet tip.}{photograph at the 14 Oct setup}

\paragraph{Motion capture.} OptiTrack tracks one rigid body per UAV and one for the ring. Its poses
are the only state measurement: the velocities and body rates of the UAVs and of the ring are
differentiated from them.
